\textbf{Proof.}
Simultaneous decomposition into \(\operatorname{span}\{\mathbf1\}\) and \(\mathbf1^\perp\) reduces \(729\Omega\) to
\[
\begin{bmatrix}60&96\\96&186\end{bmatrix},
\qquad
\begin{bmatrix}108&-108\\-108&162\end{bmatrix}.
\tag{1}
\]
Their determinants are \(1944\) and \(5832\), with positive leading entries, so \(\Omega\succ0\).  Therefore rank zero has both columns zero, while rank one has exactly one nonzero arm column.

For a focal-only normalized vector \(x\), put \(u=(\mathbf1^{\top}x)^2\).  Cauchy--Schwarz and \(\lVert x\rVert_2^2=3\) give \(0\le u\le9\).  Substitution in the scalar projector quotient from lem:rank-stratified-projector-reduction gives
\[
g_x^+(c_L)=\frac{1195u+6561}{1200(81-4u)},
\qquad
g_x^+(c_U)=\frac{243(49-5u)}{400(81-4u)}.
\tag{2}
\]
Since \(0\) is the midpoint of the two endpoints and the gain is affine,
\[
g_x^+(0)=\frac{21141-1225u}{1200(81-4u)}.
\tag{3}
\]
Using the two uncontrolled endpoint variances in def:gain-loss and taking the larger residual gives
\[
w^+(x)=\frac{9(1755-17u)}{400(81-4u)}.
\tag{4}
\]
The derivative numerator in (3) is \(-14661\), while that in (4) is \(50787\), after deletion of their positive squared denominators.  Thus (3) decreases and (4) increases on \([0,9]\), giving
\[
\sup_xg_x^+(0)=\frac{87}{400},
\qquad
\inf_xw^+(x)=\frac{39}{80},
\tag{5}
\]
both at \(u=0\).

For a comparison-only normalized vector \(y\), the same calculation gives
\[
g_y^+(c_L)=\frac{1240u+729}{150(4u+243)},
\qquad
g_y^+(c_U)=\frac{81(48-5u)}{50(4u+243)},
\tag{6}
\]
and hence
\[
g_y^+(0)=\frac{25u+12393}{300(4u+243)},
\qquad
w^+(y)=\frac{9(256u+2889)}{200(4u+243)}.
\tag{7}
\]
The derivative numerators are \(-43497\) and \(50652\), respectively.  Therefore
\[
\sup_yg_y^+(0)=\frac{17}{100},
\qquad
\inf_yw^+(y)=\frac{107}{200}.
\tag{8}
\]
At rank zero the projector is zero, so the gain is zero and the larger uncontrolled endpoint loss is \(171/200\).  Comparing (5), (8), and rank zero proves
\[
\sup_{\operatorname{rank}(S_B)\le1}g_B^+(0)=\frac{87}{400},
\qquad
\inf_{\operatorname{rank}(S_B)\le1}w^+(B)=\frac{39}{80}.
\]
Finally,
\[
\frac{31}{80}-\frac{87}{400}=\frac{17}{100}>0,
\qquad
\frac{39}{80}-\frac{2613}{8432}=\frac{468}{2635}>0,
\]
which proves both strict comparisons.  \(\square\)